% vim: tw=78 ai spell spelllang=cs
\documentclass{article}

%\usepackage{a4wide}
\usepackage{graphicx}
\usepackage[utf8x]{inputenc}
\usepackage[czech]{babel}
\usepackage[IL2]{fontenc}

\pagestyle{empty}

\begin{document}
\begin{flushright}
\today
\end{flushright}
\begin{center}
\Large\textbf{Posudek oponenta na bakalářskou práci \\
\normalsize%
\textit{Jan Dupal:} Reduced product of abstract domains}
\end{center}

Dle zadání měla bakalářská práce doimplementovat do nástroje \emph{Canal}
abstraktní interpretaci, která využije současně naimplementované interpretace
a bude je pomocí předávání informací mezi nimi využívat k přesnější analýze.
Text práce je na 29 stranách strukturovaný do 6 kapitol. V úvodu autor krátce
popisuje stav a motivaci pro práci. Následující kapitola popisuje teoretické
základy, které jsou základem pro další práci. Dále je v práci kapitola
popisující pojem ,,Reduced Product''. Obě tyto kapitoly jsou podkladem pro
následující kapitolu o implementaci. Zbývá kapitola o porovnání s původní
implementací. Práci ukončuje závěr obsahující i krátký popis práce do
budoucna.  \textbf{Rozsah} uvedených kapitol je z pohledu bakalářské práce i
zadání obvyklý.

\textbf{Obsahově} dle mého názoru dostatečně popisuje teoretické základy
práce, implementaci i kvantifikuje její přínos. Je tu zde jen několik
nesrovnalostí (např.  v bodu 1 definice 2 má být jeden ze symbolů $\gamma$).

Po \textbf{jazykové} stránce je práce na velmi slušné úrovni. Je zde jen
minimum překlepů a chyb (např. ,,this abstractions'' na str. 2 nebo ,,products
performs'' na str. 15). Oceňuji také volbu anglického jazyka.

\textbf{Praktická část} práce je v přiloženém \texttt{zip} souboru. Jelikož se
jedná o rozšíření již existujícího nástroje \emph{Canal}, je v práci přesně
uvedeno, co přesně bylo změněno a přidáno. Tento nově přidaný kód je čitelný a
jeví se být funkční.

Na závěr konstatuji, že práce \textbf{splňuje zadání a požadavky na
bakalářskou práci}. Navrhuji práci \textbf{uznat} jako bakalářskou a vzhledem
k výše uvedeným okolnostem doporučuji hodnocení stupněm A.

\vspace{3em}
\begin{flushright}
Jiří Slabý
\end{flushright}

\end{document}
